\chapter{预备知识}
\label{ch:prelim}

\begin{flushright}
\textit{“引其纲，万目皆张。”——《吕氏春秋·用民》}
\footnote{纲是网上的总绳.这一章的纲是矢量、梯度和指标记号,抓住以后,散度定理、正交坐标和张量恒等式都顺着张开.}
\end{flushright}

后面几章用到矢量、线性映射和正交曲线坐标.这一章把它们按计算时需要的顺序写在一起:先是矢量的乘积,再是矩阵的逆,然后在直角坐标和正交曲线坐标里写出梯度、散度和旋度,最后用指标记号和张量把恒等式收拢.同一套材料的标准讨论见~\cite{arfken,byron,riley}.

\section{矢量}
\idx{矢量}{shiliang}

标量用一个数就够了,坐标旋转时它不变.矢量既有大小也有方向.在三维直角坐标里把它写成
\begin{equation}
  \mathbf{A}=A_x\mathbf{e}_x+A_y\mathbf{e}_y+A_z\mathbf{e}_z,
\end{equation}
基矢满足 $\mathbf{e}_i\cdot\mathbf{e}_j=\delta_{ij}$.模是
\begin{equation}
  |\mathbf{A}|=\sqrt{A_x^2+A_y^2+A_z^2}.
\end{equation}
数乘 $c\mathbf{A}$ 把每个分量乘上 $c$.$c>0$ 时方向不变,$c<0$ 时方向反过来.

\subsection{加法与减法}
分量各自相加.几何上,$\mathbf{A}+\mathbf{B}$ 是以 $\mathbf{A},\mathbf{B}$ 为邻边的平行四边形的对角线;$\mathbf{B}-\mathbf{A}$ 是从 $\mathbf{A}$ 的终点指向 $\mathbf{B}$ 的终点.

\begin{figure}[htb]
  \centering
  \begin{tikzpicture}[
    font=\small, >=Stealth,
    vec/.style={-{Stealth[length=2mm]}, line width=0.8pt},
  ]
  \def\s{0.28}
  \begin{scope}
    \coordinate (O) at (0,0);
    \coordinate (A) at (4*\s,4*\s);
    \coordinate (B) at (1*\s,8*\s);
    \coordinate (S) at (5*\s,12*\s);
    \fill[blue!8] (O) -- (A) -- (S) -- (B) -- cycle;
    \draw[gray!45] (A) -- (S) -- (B);
    \draw[vec, blue!70!black] (O) -- (A) node[left] {$\mathbf{A}$};
    \draw[vec, red!70!black] (O) -- (B) node[left] {$\mathbf{B}$};
    \draw[vec, purple!70!black] (O) -- (S) node[right] {$\mathbf{A}+\mathbf{B}$};
    \node at (0.7,-0.55) {（a）平行四边形};
  \end{scope}
  \begin{scope}[xshift=4.6cm]
    \coordinate (O) at (0,0);
    \coordinate (A) at (4*\s,4*\s);
    \coordinate (B) at (1*\s,8*\s);
    \draw[vec, blue!70!black] (O) -- (A) node[left] {$\mathbf{A}$};
    \draw[vec, red!70!black] (O) -- (B) node[right] {$\mathbf{B}$};
    \draw[vec, purple!70!black] (A) -- (B) node[midway, left] {$\mathbf{B}-\mathbf{A}$};
    \node at (0.55,-0.55) {（b）$\mathbf{B}-\mathbf{A}$};
  \end{scope}
\end{tikzpicture}

  \caption{平行四边形给出和;把两个矢量的起点放在一起,终点的连线是差.}
  \label{fig:sum-diff}
\end{figure}

\begin{example}
$\mathbf{A}=4\mathbf{e}_x+4\mathbf{e}_y$,$\mathbf{B}=\mathbf{e}_x+8\mathbf{e}_y$.
\end{example}
\begin{solution}
\begin{equation}
  \mathbf{A}+\mathbf{B}=5\mathbf{e}_x+12\mathbf{e}_y,
  \qquad
  |\mathbf{A}+\mathbf{B}|=\sqrt{5^2+12^2}=13.
\end{equation}
\begin{equation}
  \mathbf{B}-\mathbf{A}=-3\mathbf{e}_x+4\mathbf{e}_y,
  \qquad
  |\mathbf{B}-\mathbf{A}|=\sqrt{(-3)^2+4^2}=5.
\end{equation}
$5$--$12$--$13$ 和 $3$--$4$--$5$ 两组勾股数正好核对这两个模.
\end{solution}

\subsection{点积}
\idx{点积}{dianji}
\begin{equation}
  \mathbf{A}\cdot\mathbf{B}=|\mathbf{A}|\,|\mathbf{B}|\cos\theta
  =A_x B_x+A_y B_y+A_z B_z.
\end{equation}
$\theta$ 是夹角.点积为正是锐角,为零是垂直,为负是钝角.特别地 $\mathbf{A}\cdot\mathbf{A}=|\mathbf{A}|^2$.

\begin{example}
$\mathbf{A}=\sqrt{3}\,\mathbf{e}_x+\mathbf{e}_y$,$\mathbf{B}=2\mathbf{e}_x$.
\end{example}
\begin{solution}
$|\mathbf{A}|=2$,$|\mathbf{B}|=2$,$\mathbf{A}\cdot\mathbf{B}=2\sqrt{3}$.因此 $\cos\theta=\sqrt{3}/2$,$\theta=30^\circ$.$\mathbf{B}$ 沿 $x$ 轴,$\mathbf{A}$ 与它成 $30^\circ$.
\end{solution}

\begin{example}
水分子的键角.氧原子在原点,两个氢原子的坐标(单位 \AA)是 $\mathrm{H}_1=(0.757,0,0.586)$,$\mathrm{H}_2=(-0.757,0,0.586)$.
\end{example}
\begin{solution}
键矢量 $\mathbf{u}=(a,0,b)$,$\mathbf{v}=(-a,0,b)$,$a=0.757$,$b=0.586$.
\begin{equation}
  \cos\alpha=\frac{\mathbf{u}\cdot\mathbf{v}}{|\mathbf{u}|\,|\mathbf{v}|}
  =\frac{b^2-a^2}{a^2+b^2}.
\end{equation}
$a^2=0.573$,$b^2=0.343$,所以 $\cos\alpha\approx -0.251$,$\alpha\approx 104.5^\circ$.
\end{solution}

\begin{example}
$\mathbf{p}=2\mathbf{e}_x-\mathbf{e}_y+2\mathbf{e}_z$ 与 $\mathbf{q}=\mathbf{e}_x+2\mathbf{e}_y$ 是否垂直.
\end{example}
\begin{solution}
$\mathbf{p}\cdot\mathbf{q}=2-2+0=0$,所以垂直.
\end{solution}

\subsection{叉积}
\idx{叉积}{chaji}
\begin{equation}
  \mathbf{A}\times\mathbf{B}
  =
  \begin{vmatrix}
    \mathbf{e}_x & \mathbf{e}_y & \mathbf{e}_z \\
    A_x & A_y & A_z \\
    B_x & B_y & B_z
  \end{vmatrix}.
\end{equation}
它垂直于 $\mathbf{A}$ 和 $\mathbf{B}$,方向由右手定则决定,模 $|\mathbf{A}|\,|\mathbf{B}|\sin\theta$ 是二者张成的平行四边形面积.$\mathbf{A}\times\mathbf{A}=\mathbf{0}$,并且 $\mathbf{B}\times\mathbf{A}=-\mathbf{A}\times\mathbf{B}$.

\begin{example}
$\mathbf{A}=-\mathbf{e}_x+\mathbf{e}_y+\mathbf{e}_z$,$\mathbf{B}=\mathbf{e}_x-\mathbf{e}_y+\mathbf{e}_z$.
\end{example}
\begin{solution}
按第一行展开,
\begin{equation}
  \mathbf{A}\times\mathbf{B}=2\mathbf{e}_x+2\mathbf{e}_y,
  \qquad
  |\mathbf{A}\times\mathbf{B}|=2\sqrt{2}.
\end{equation}
右手单位法向是 $(\mathbf{e}_x+\mathbf{e}_y)/\sqrt{2}$.同一对矢量的点积是 $-1$,两个模都是 $\sqrt{3}$,所以
\begin{equation}
  \cos\theta=-\frac13,
  \qquad
  \theta=\cos^{-1}(-1/3)\approx 109.47^\circ.
\end{equation}
叉积给出法向,点积给出平面内的夹角.
\end{solution}

\subsection{三重积}
\idx{三重积}{sanzhongji}
标量三重积
\begin{equation}
  \mathbf{A}\cdot(\mathbf{B}\times\mathbf{C})
  =
  \begin{vmatrix}
    A_x & A_y & A_z \\
    B_x & B_y & B_z \\
    C_x & C_y & C_z
  \end{vmatrix}
\end{equation}
的绝对值是三个矢量所张平行六面体的体积,符号记录定向.把 $(\mathbf{A}\times\mathbf{B})\cdot\mathbf{C}$ 展开后与上式相同,所以点与叉可以交换.记 $[\mathbf{A},\mathbf{B},\mathbf{C}]=\mathbf{A}\cdot(\mathbf{B}\times\mathbf{C})$,循环换位不变号,交换两个矢量变号.

\begin{example}
$\mathbf{A}=\mathbf{e}_x+\mathbf{e}_y$,$\mathbf{B}=\mathbf{e}_y+\mathbf{e}_z$,$\mathbf{C}=\mathbf{e}_x+\mathbf{e}_z$.
\end{example}
\begin{solution}
$\mathbf{A}\times\mathbf{B}=\mathbf{e}_x-\mathbf{e}_y+\mathbf{e}_z$,再点乘 $\mathbf{C}$ 得 $2$.行列式
\begin{equation}
  \begin{vmatrix}
    1 & 1 & 0 \\
    0 & 1 & 1 \\
    1 & 0 & 1
  \end{vmatrix}
  =2
\end{equation}
给出同一个有向体积.
\end{solution}

两个叉积的点积是拉格朗日恒等式
\begin{equation}
  (\mathbf{A}\times\mathbf{B})\cdot(\mathbf{C}\times\mathbf{D})
  =(\mathbf{A}\cdot\mathbf{C})(\mathbf{B}\cdot\mathbf{D})
  -(\mathbf{A}\cdot\mathbf{D})(\mathbf{B}\cdot\mathbf{C}).
  \label{eq:lagrange-identity}
\end{equation}
用 $\varepsilon_{ijk}\varepsilon_{ilm}=\delta_{jl}\delta_{km}-\delta_{jm}\delta_{kl}$ 把指标收成点积就得到它.这个缩并放在指标记号一节里证明.$|\mathbf{A}\times\mathbf{B}|^2$ 是 $\mathbf{C}=\mathbf{A}$,$\mathbf{D}=\mathbf{B}$ 的特例.

矢量三重积没有结合律.化简公式是
\begin{equation}
\begin{aligned}
  \mathbf{A}\times(\mathbf{B}\times\mathbf{C})
  &=(\mathbf{A}\cdot\mathbf{C})\mathbf{B}-(\mathbf{A}\cdot\mathbf{B})\mathbf{C}, \\
  (\mathbf{A}\times\mathbf{B})\times\mathbf{C}
  &=(\mathbf{A}\cdot\mathbf{C})\mathbf{B}-(\mathbf{B}\cdot\mathbf{C})\mathbf{A}.
\end{aligned}
\label{eq:bac-cab}
\end{equation}
第一条的 $x$ 分量可以直接展开:先算 $(\mathbf{B}\times\mathbf{C})_y$ 和 $(\mathbf{B}\times\mathbf{C})_z$,再与 $A_y,A_z$ 相乘,整理后就是 $(\mathbf{A}\cdot\mathbf{C})B_x-(\mathbf{A}\cdot\mathbf{B})C_x$.另外两个分量相同.把三个矢量循环置换,得到第二条.三条公式相加,$\mathbf{A},\mathbf{B},\mathbf{C}$ 的系数各自抵消,所以
\begin{equation}
  \mathbf{A}\times(\mathbf{B}\times\mathbf{C})
  +\mathbf{B}\times(\mathbf{C}\times\mathbf{A})
  +\mathbf{C}\times(\mathbf{A}\times\mathbf{B})
  =\mathbf{0}.
\end{equation}

\begin{example}
单位矢量 $\hat{\mathbf{n}}$ 与 $\mathbf{v}=2\mathbf{e}_x-\mathbf{e}_y+3\mathbf{e}_z$,取 $\hat{\mathbf{n}}=\mathbf{e}_z$.
\end{example}
\begin{solution}
由\eqref{eq:bac-cab},
\begin{equation}
  \hat{\mathbf{n}}\times(\mathbf{v}\times\hat{\mathbf{n}})
  =\mathbf{v}-(\mathbf{v}\cdot\hat{\mathbf{n}})\hat{\mathbf{n}}.
\end{equation}
这里 $\mathbf{v}\cdot\mathbf{e}_z=3$,垂直分量是 $2\mathbf{e}_x-\mathbf{e}_y$.减去的那一项是沿 $\hat{\mathbf{n}}$ 的投影.
\end{solution}

\section{矩阵与逆}
\idx{矩阵}{juzhen}
\idx{逆矩阵}{nijuzhen}

把矢量写成列.矩阵 $M$ 作用在 $\mathbf{x}$ 上,是各列的线性组合,
\begin{equation}
  (M\mathbf{x})_i=M_{ij}x_j.
\end{equation}
重复指标 $j$ 要求和,这是下一节的求和约定,这里先按这个意思使用.相继两次映射先做 $N$ 再做 $M$,
\begin{equation}
  M(N\mathbf{x})=(MN)\mathbf{x}.
\end{equation}
单位矩阵的列就是基矢,$I\mathbf{x}=\mathbf{x}$.单位正方形被 $M$ 送到一个平行四边形,面积是 $|\det M|$.

\subsection{逆矩阵}
若存在 $A^{-1}$ 使 $A A^{-1}=A^{-1}A=I$,就说 $A$ 可逆,线性映射可以做回去.逆存在当且仅当 $\det A\neq 0$.用余子式构造:划去第 $i$ 行第 $j$ 列,剩下的行列式乘上 $(-1)^{i+j}$,得到余子式 $C_{ij}$.伴随矩阵是 $\operatorname{adj} A=C^{\mathsf T}$,并且
\begin{equation}
  A\,(\operatorname{adj} A)=(\det A)\,I,
  \qquad
  A^{-1}=\frac{1}{\det A}\operatorname{adj} A.
\end{equation}

$2\times 2$ 可以直接验证.设
\begin{equation}
  A=\begin{pmatrix} a & b \\ c & d \end{pmatrix},
  \qquad
  \det A=ad-bc.
\end{equation}
\begin{equation}
  \begin{pmatrix} a & b \\ c & d \end{pmatrix}
  \begin{pmatrix} d & -b \\ -c & a \end{pmatrix}
  =(ad-bc)I,
\end{equation}
所以
\begin{equation}
  A^{-1}=\frac{1}{ad-bc}\begin{pmatrix} d & -b \\ -c & a \end{pmatrix}.
\end{equation}

\begin{example}
$A=\begin{pmatrix}2&1\\1&2\end{pmatrix}$,$\mathbf{b}=\begin{pmatrix}1\\0\end{pmatrix}$.
\end{example}
\begin{solution}
$\det A=3$,
\begin{equation}
  A^{-1}=\frac13\begin{pmatrix}2&-1\\-1&2\end{pmatrix}.
\end{equation}
直接相乘,$AA^{-1}=I$.方程 $A\mathbf{x}=\mathbf{b}$ 的解是
\begin{equation}
  \mathbf{x}=A^{-1}\mathbf{b}=\begin{pmatrix}2/3\\-1/3\end{pmatrix}.
\end{equation}
\end{solution}

\begin{example}
\begin{equation}
  A=\begin{pmatrix} 2&1&0\\ 1&2&1\\ 0&1&2 \end{pmatrix}.
\end{equation}
\end{example}
\begin{solution}
沿第一行展开,$\det A=2(4-1)-(2-0)=4$.余子式矩阵是
\begin{equation}
  C=\begin{pmatrix} 3&-2&1\\ -2&4&-2\\ 1&-2&3 \end{pmatrix}.
\end{equation}
它已经对称,$\operatorname{adj} A=C$,
\begin{equation}
  A^{-1}=\frac14 C.
\end{equation}
$A C$ 的 $(1,1)$ 元是 $2\cdot 3+1\cdot(-2)=4=\det A$,$(1,2)$ 元是 $2\cdot(-2)+1\cdot 4=0$.这与 $A\,(\operatorname{adj} A)=(\det A)\,I$ 一致.
\end{solution}

若 $R^{\mathsf T}R=I$,则 $R^{-1}=R^{\mathsf T}$,不必再做余子式.旋转矩阵
\begin{equation}
  R=\begin{pmatrix} \cos\theta & -\sin\theta \\ \sin\theta & \cos\theta \end{pmatrix}
\end{equation}
就是这种正交矩阵.它保持矢量的模和夹角,核对放在指标记号里.

\section{梯度、散度与旋度}
\idx{梯度}{tidu}
\idx{散度}{sandu}
\idx{旋度}{xuandu}
\label{sec:coordinates}

这三个算子由弧长、通量和环量定义.直角坐标的尺度因子都是 $1$.柱坐标和球坐标的弧长不是坐标增量本身,同一套定义要带上尺度因子.矩形边界用直角坐标,圆和圆柱用柱坐标,球用球坐标;后面在这些边界上分离变量时,直接使用本节的公式.

\subsection{尺度因子}
\idx{尺度因子}{chiduyinzi}
\idx{正交曲线坐标}{zhengjiaoxianzuobiao}

直角坐标系 $(x,y,z)$、柱坐标系 $(\rho,\phi,z)$ 和球坐标系 $(r,\theta,\phi)$ 都是正交曲线坐标系:三族坐标面彼此正交.一般的曲线坐标记作 $(q_1,q_2,q_3)$,与直角坐标的变换
\begin{equation}
  x=x(q_1,q_2,q_3),\qquad
  y=y(q_1,q_2,q_3),\qquad
  z=z(q_1,q_2,q_3)
\end{equation}
在所讨论的区域里一一对应,雅可比行列式不为零,
\begin{equation}
  \frac{\partial(x,y,z)}{\partial(q_1,q_2,q_3)}\neq 0.
\end{equation}

位置矢量 $\br=\br(q_1,q_2,q_3)$.沿第 $i$ 条坐标线只变动 $q_i$,切向量是 $\partial\br/\partial q_i$.若不同坐标线的切向量互相正交,
\begin{equation}
  \frac{\partial\br}{\partial q_i}\cdot\frac{\partial\br}{\partial q_j}=0,
  \qquad i\neq j,
\end{equation}
这个坐标系就是正交的.弧长元没有交叉项,
\begin{equation}
  \mathrm{d}s^2
  =\mathrm{d}\br\cdot\mathrm{d}\br
  =\sum_{i=1}^{3} h_i^2\,\mathrm{d}q_i^2,
  \qquad
  h_i=\left|\frac{\partial\br}{\partial q_i}\right|.
  \label{eq:scale-factor}
\end{equation}
$h_i$ 称为度规系数,或尺度因子.它把坐标增量 $\mathrm{d}q_i$ 换成真正的长度 $h_i\,\mathrm{d}q_i$.直角坐标的三个尺度因子都是 $1$.

第 $i$ 个单位矢量沿坐标线的切线,
\begin{equation}
  \mathbf{e}_i=\frac{1}{h_i}\frac{\partial\br}{\partial q_i}.
  \label{eq:unit-vector}
\end{equation}
正交性写成 $\mathbf{e}_i\cdot\mathbf{e}_j=\delta_{ij}$.和直角坐标不同,$\mathbf{e}_i$ 一般随点改变.算散度和旋度时,不能把它们当成常矢量提出求导.

\begin{figure}[htb]
  \centering
  \import{curvedcoordinates/tikz/}{bases.tex}
  \caption{柱坐标和球坐标的局部单位矢量.柱坐标的 $\mathbf{e}_z$ 垂直于纸面;球坐标画的是过极轴的子午面,$\mathbf{e}_{\phi}$ 指向纸外.}
  \label{fig:curvilinear-bases}
\end{figure}

\begin{example}
由位置矢量求柱坐标和球坐标的尺度因子与单位矢量.
\end{example}
\begin{solution}
柱坐标 $\br=\rho\cos\phi\,\mathbf{e}_x+\rho\sin\phi\,\mathbf{e}_y+z\,\mathbf{e}_z$.
\begin{align}
  \frac{\partial\br}{\partial\rho}
  &=\cos\phi\,\mathbf{e}_x+\sin\phi\,\mathbf{e}_y, \\
  \frac{\partial\br}{\partial\phi}
  &=-\rho\sin\phi\,\mathbf{e}_x+\rho\cos\phi\,\mathbf{e}_y, \\
  \frac{\partial\br}{\partial z}
  &=\mathbf{e}_z.
\end{align}
三式两两正交,并且
\begin{equation}
  h_{\rho}=1,\qquad h_{\phi}=\rho,\qquad h_z=1.
\end{equation}
$\phi$ 增加 $\mathrm{d}\phi$ 时,走过的弧长是 $\rho\,\mathrm{d}\phi$,不是 $\mathrm{d}\phi$.单位矢量是
\begin{equation}
  \mathbf{e}_{\rho}=\cos\phi\,\mathbf{e}_x+\sin\phi\,\mathbf{e}_y,
  \qquad
  \mathbf{e}_{\phi}=-\sin\phi\,\mathbf{e}_x+\cos\phi\,\mathbf{e}_y,
  \qquad
  \mathbf{e}_z=\mathbf{e}_z.
\end{equation}

球坐标 $\br=r\sin\theta\cos\phi\,\mathbf{e}_x+r\sin\theta\sin\phi\,\mathbf{e}_y+r\cos\theta\,\mathbf{e}_z$.
\begin{align}
  \frac{\partial\br}{\partial r}
  &=\sin\theta\cos\phi\,\mathbf{e}_x+\sin\theta\sin\phi\,\mathbf{e}_y+\cos\theta\,\mathbf{e}_z, \\
  \frac{\partial\br}{\partial\theta}
  &=r(\cos\theta\cos\phi\,\mathbf{e}_x+\cos\theta\sin\phi\,\mathbf{e}_y-\sin\theta\,\mathbf{e}_z), \\
  \frac{\partial\br}{\partial\phi}
  &=r\sin\theta(-\sin\phi\,\mathbf{e}_x+\cos\phi\,\mathbf{e}_y).
\end{align}
同样两两正交,
\begin{equation}
  h_r=1,\qquad h_{\theta}=r,\qquad h_{\phi}=r\sin\theta.
\end{equation}
单位矢量就是上面三个偏导各自除以对应的 $h_i$.纬度圈上的弧长是 $r\sin\theta\,\mathrm{d}\phi$,所以 $\phi$ 方向的尺度因子在两极变为零.
\end{solution}

\begin{figure}[htb]
  \centering
  \import{curvedcoordinates/tikz/}{volume.tex}
  \caption{正交曲线坐标中的体积元.三条棱长分别是 $h_i\,\mathrm{d}q_i$,体积是三者之积.}
  \label{fig:volume-element}
\end{figure}

沿坐标线的长度、坐标面上的面积和体积元是
\begin{align}
  \mathrm{d}\ell_i &= h_i\,\mathrm{d}q_i, \\
  \mathrm{d}S_1 &= h_2 h_3\,\mathrm{d}q_2\,\mathrm{d}q_3, \\
  \mathrm{d}V &= h_1 h_2 h_3\,\mathrm{d}q_1\,\mathrm{d}q_2\,\mathrm{d}q_3.
  \label{eq:volume-element}
\end{align}
$\mathrm{d}S_1$ 是法向为 $\mathbf{e}_1$ 的那一面.柱坐标 $\mathrm{d}V=\rho\,\mathrm{d}\rho\,\mathrm{d}\phi\,\mathrm{d}z$,球坐标 $\mathrm{d}V=r^2\sin\theta\,\mathrm{d}r\,\mathrm{d}\theta\,\mathrm{d}\phi$.雅可比行列式的绝对值等于 $h_1 h_2 h_3$.

\subsection{梯度}

标量场 $u$ 沿位移 $\mathrm{d}\boldsymbol{\ell}$ 的变化是 $\mathrm{d}u=\nabla u\cdot\mathrm{d}\boldsymbol{\ell}$.在局部正交归一基上,$\nabla u=\sum f_i\mathbf{e}_i$,方向导数等于沿弧长的变化率,
\begin{equation}
  f_i=\mathbf{e}_i\cdot\nabla u
  =\frac{\partial u}{\partial\ell_i}
  =\frac{1}{h_i}\frac{\partial u}{\partial q_i}.
\end{equation}
因此
\begin{equation}
  \nabla u
  =\sum_{i=1}^{3}\frac{1}{h_i}\frac{\partial u}{\partial q_i}\,\mathbf{e}_i.
  \label{eq:grad-orthogonal}
\end{equation}
直角坐标里 $h_i=1$,$\mathrm{d}\ell_i$ 就是 $\mathrm{d}x,\mathrm{d}y,\mathrm{d}z$,全微分
\begin{equation}
  \mathrm{d}u
  =\frac{\partial u}{\partial x}\,\mathrm{d}x
  +\frac{\partial u}{\partial y}\,\mathrm{d}y
  +\frac{\partial u}{\partial z}\,\mathrm{d}z
\end{equation}
给出
\begin{equation}
  \nabla u
  =\frac{\partial u}{\partial x}\,\mathbf{e}_x
  +\frac{\partial u}{\partial y}\,\mathbf{e}_y
  +\frac{\partial u}{\partial z}\,\mathbf{e}_z.
\end{equation}
$\nabla u$ 指向 $u$ 增加最快的方向,它的模是这个方向上的方向导数.等值面 $u=\textrm{常数}$ 的切向位移使 $\mathrm{d}u=0$,所以梯度垂直于等值面.尺度因子一旦不是 $1$,相应的偏导前面就要除以它.

柱坐标和球坐标分别是
\begin{align}
  \nabla u
  &=\frac{\partial u}{\partial\rho}\,\mathbf{e}_{\rho}
    +\frac{1}{\rho}\frac{\partial u}{\partial\phi}\,\mathbf{e}_{\phi}
    +\frac{\partial u}{\partial z}\,\mathbf{e}_z,
  \label{eq:grad-cyl} \\
  \nabla u
  &=\frac{\partial u}{\partial r}\,\mathbf{e}_r
    +\frac{1}{r}\frac{\partial u}{\partial\theta}\,\mathbf{e}_{\theta}
    +\frac{1}{r\sin\theta}\frac{\partial u}{\partial\phi}\,\mathbf{e}_{\phi}.
  \label{eq:grad-sph}
\end{align}

沿路径的功 $\int_L\mathbf{F}\cdot\mathrm{d}\boldsymbol{\ell}$,在 $\mathbf{F}=\nabla u$ 时只依赖端点,
\begin{equation}
  \int_A^B\nabla u\cdot\mathrm{d}\boldsymbol{\ell}=u(B)-u(A).
\end{equation}
闭路上积分是零.

\begin{figure}[htb]
  \centering
  \begin{tikzpicture}[font=\small, >=Stealth,
    vec/.style={-{Stealth[length=1.8mm]}, line width=0.7pt}]
  \coordinate (O) at (0,0);
  \coordinate (P) at (2.6,1.7);
  \coordinate (X) at (2.6,0);
  \coordinate (Y) at (0,1.7);
  \draw[gray!50] (-0.2,0) -- (3.1,0) node[right] {$x$};
  \draw[gray!50] (0,-0.2) -- (0,2.15) node[above] {$y$};
  \draw[vec, blue!70!black] (O) -- (X) -- (P);
  \draw[vec, red!70!black] (O) -- (Y) -- (P);
  \fill (O) circle (1.1pt) node[below left] {$O$};
  \fill (P) circle (1.1pt) node[above right] {$P$};
  \node[blue!70!black, font=\scriptsize] at (1.5,-0.35) {$L_1$};
  \node[red!70!black, font=\scriptsize] at (-0.35,0.85) {$L_2$};
\end{tikzpicture}

  \caption{从原点到 $P$ 的两条折线.场如果是梯度,两条路径的积分相等.}
  \label{fig:two-paths}
\end{figure}

\begin{example}
$u=x^2 y$,从原点到 $P=(x_0,y_0)$.
\end{example}
\begin{solution}
$\nabla u=2xy\,\mathbf{e}_x+x^2\,\mathbf{e}_y$,端点公式给出 $x_0^2 y_0$.

路径 $L_1$ 先沿 $x$ 轴到 $(x_0,0)$,再竖直到 $P$.第一段 $y=0$,被积函数是零;第二段 $x=x_0$,$\int_0^{y_0} x_0^2\,\mathrm{d}y=x_0^2 y_0$.

路径 $L_2$ 先沿 $y$ 轴到 $(0,y_0)$,再水平到 $P$.第一段 $x=0$,被积函数是零;第二段 $\int_0^{x_0} 2x y_0\,\mathrm{d}x=x_0^2 y_0$.

两条路径结果相同.任何其他路径也给出这个值.
\end{solution}

\begin{example}
$u=\rho^2$ 与 $u=\ln\rho$ 的梯度.
\end{example}
\begin{solution}
$u$ 只依赖 $\rho$.由\eqref{eq:grad-cyl},
\begin{equation}
  \nabla(\rho^2)=2\rho\,\mathbf{e}_{\rho},
  \qquad
  \nabla\ln\rho=\frac{1}{\rho}\,\mathbf{e}_{\rho}
  \quad(\rho>0).
\end{equation}
第二式的长度是 $1/\rho$,方向径向向外.它就是二维点源电场的形状.
\end{solution}

\subsection{散度}

闭曲面的通量是 $\oint_S\mathbf{A}\cdot\mathrm{d}\mathbf{S}$,$\mathrm{d}\mathbf{S}$ 取外法向.净通量为正,内部有源.把矢量按同一组单位矢量分解,$\mathbf{A}=\sum A_i\mathbf{e}_i$,$A_i=\mathbf{A}\cdot\mathbf{e}_i$ 是物理分量.散度是体积元表面的净通量除以体积.

考虑图\ref{fig:volume-element} 那个小盒子.法向为 $\mathbf{e}_1$ 的两面,面积是 $h_2 h_3\,\mathrm{d}q_2\,\mathrm{d}q_3$,通量之差是
\begin{equation}
  \frac{\partial}{\partial q_1}(A_1 h_2 h_3)\,\mathrm{d}q_1\,\mathrm{d}q_2\,\mathrm{d}q_3.
\end{equation}
三个方向相加,再除以 $\mathrm{d}V=h_1 h_2 h_3\,\mathrm{d}q_1\,\mathrm{d}q_2\,\mathrm{d}q_3$,
\begin{equation}
  \nabla\cdot\mathbf{A}
  =\frac{1}{h_1 h_2 h_3}
    \sum_{i=1}^{3}
    \frac{\partial}{\partial q_i}
    \Bigl(\frac{h_1 h_2 h_3}{h_i} A_i\Bigr).
  \label{eq:div-orthogonal}
\end{equation}
尺度因子留在导数里面.如果先把 $h_i$ 提出来,就会丢掉坐标基元本身的变化.直角坐标里 $h_i=1$,上式回到
\begin{equation}
  \nabla\cdot\mathbf{A}
  =\frac{\partial A_x}{\partial x}
  +\frac{\partial A_y}{\partial y}
  +\frac{\partial A_z}{\partial z}.
\end{equation}
边长 $\Delta x,\Delta y,\Delta z$ 的小长方体上,在 $x$ 与 $x+\Delta x$ 这一对面上的净通量就是 $(\partial A_x/\partial x)\,\Delta V$.

内部小面的通量成对抵消,只留下外表面,这就是散度定理
\idx{散度定理}{sandudingli}
\begin{equation}
  \oint_S\mathbf{A}\cdot\mathrm{d}\mathbf{S}
  =\int_V(\nabla\cdot\mathbf{A})\,\mathrm{d}V.
  \label{eq:div-theorem}
\end{equation}

柱坐标 $h_{\rho}h_{\phi}h_z=\rho$,
\begin{equation}
  \nabla\cdot\mathbf{A}
  =\frac{1}{\rho}\frac{\partial}{\partial\rho}(\rho A_{\rho})
   +\frac{1}{\rho}\frac{\partial A_{\phi}}{\partial\phi}
   +\frac{\partial A_z}{\partial z}.
  \label{eq:div-cyl}
\end{equation}
球坐标 $h_r h_{\theta} h_{\phi}=r^2\sin\theta$,
\begin{equation}
  \nabla\cdot\mathbf{A}
  =\frac{1}{r^2}\frac{\partial}{\partial r}(r^2 A_r)
   +\frac{1}{r\sin\theta}\frac{\partial}{\partial\theta}(\sin\theta\, A_{\theta})
   +\frac{1}{r\sin\theta}\frac{\partial A_{\phi}}{\partial\phi}.
  \label{eq:div-sph}
\end{equation}

\begin{example}
$\mathbf{A}=x\mathbf{e}_x+y\mathbf{e}_y+z\mathbf{e}_z$,长方体 $0\leq x\leq a$,$0\leq y\leq b$,$0\leq z\leq c$.
\end{example}
\begin{solution}
$\nabla\cdot\mathbf{A}=3$,体积分是 $3abc$.六个面里,过原点的三个面通量是零,对面分别贡献 $abc$.合计仍是 $3abc$.
\end{solution}

\begin{example}
$\mathbf{A}=\mathbf{e}_{\rho}$ 与 $\mathbf{A}=\mathbf{e}_r/r^2$ 的散度.
\end{example}
\begin{solution}
柱坐标里 $A_{\rho}=1$,其余分量为零.\eqref{eq:div-cyl}只留下
\begin{equation}
  \nabla\cdot\mathbf{e}_{\rho}=\frac{1}{\rho}\frac{\partial}{\partial\rho}(\rho)=\frac{1}{\rho}.
\end{equation}
三个直角分量都是常数的矢量,散度才是零.$\mathbf{e}_{\rho}$ 的方向在空间里转,散度不是零.

球坐标里 $A_r=1/r^2$.\eqref{eq:div-sph}的第一项是
\begin{equation}
  \frac{1}{r^2}\frac{\partial}{\partial r}(r^2\cdot r^{-2})=0,
  \qquad r>0.
\end{equation}
所以 $\nabla\cdot(\mathbf{e}_r/r^2)=0$（$r\neq 0$）.原点是奇点,把原点包进去以后通量是 $4\pi$,这就是 $-\Delta(1/r)=4\pi\delta^3(\br)$ 的来源.
\end{solution}

\subsection{旋度}

闭路环量是 $\oint_L\mathbf{A}\cdot\mathrm{d}\boldsymbol{\ell}$.在法向为 $\mathbf{e}_1$ 的小面片上用斯托克斯定理的局部形式.环量来自两条 $q_2$ 边和两条 $q_3$ 边,切向长度分别是 $h_2\,\mathrm{d}q_2$ 和 $h_3\,\mathrm{d}q_3$.除以面积 $h_2 h_3\,\mathrm{d}q_2\,\mathrm{d}q_3$,
\begin{equation}
  (\nabla\times\mathbf{A})_1
  =\frac{1}{h_2 h_3}
    \Biggl(
      \frac{\partial}{\partial q_2}(h_3 A_3)
      -\frac{\partial}{\partial q_3}(h_2 A_2)
    \Biggr).
  \label{eq:curl-orthogonal}
\end{equation}
另外两个分量按 $1\to 2\to 3\to 1$ 轮换.导数作用在 $h_i A_i$ 上,而不是只作用在 $A_i$ 上.直角坐标里 $h_i=1$,$xy$ 平面的小矩形给出
\begin{equation}
  (\nabla\times\mathbf{A})_z
  =\frac{\partial A_y}{\partial x}-\frac{\partial A_x}{\partial y}.
\end{equation}

\begin{figure}[htb]
  \centering
  \begin{tikzpicture}[font=\small, >=Stealth,
    vec/.style={-{Stealth[length=1.6mm]}, line width=0.7pt}]
  \coordinate (O) at (0,0);
  \coordinate (X) at (2.2,0);
  \coordinate (Y) at (0,1.5);
  \coordinate (P) at (2.2,1.5);
  \draw[vec, blue!70!black] (O) -- (X) node[midway, below] {$A_x\Delta x$};
  \draw[vec, blue!70!black] (P) -- (Y) node[midway, above] {$-A_x\Delta x$};
  \draw[vec, red!70!black] (X) -- (P) node[midway, right] {$A_y\Delta y$};
  \draw[vec, red!70!black] (Y) -- (O) node[midway, left] {$-A_y\Delta y$};
  \node at (1.1,0.75) {$(\nabla\times\mathbf{A})_z$};
\end{tikzpicture}

  \caption{小矩形上的环量.内部公共边方向相反,铺满曲面以后互相抵消.}
  \label{fig:curl-loop}
\end{figure}

一般地,法向为 $\hat{\mathbf{n}}$ 的小面元满足
\begin{equation}
  (\nabla\times\mathbf{A})\cdot\hat{\mathbf{n}}
  =\lim_{\Delta S\to 0}\frac{1}{\Delta S}
    \oint_{\partial(\Delta S)}\mathbf{A}\cdot\mathrm{d}\boldsymbol{\ell},
\end{equation}
环路与法向由右手定则联系.内部边抵消后就是斯托克斯定理
\idx{斯托克斯定理}{situokesidingli}
\begin{equation}
  \oint_L\mathbf{A}\cdot\mathrm{d}\boldsymbol{\ell}
  =\int_S(\nabla\times\mathbf{A})\cdot\mathrm{d}\mathbf{S}.
  \label{eq:stokes}
\end{equation}
张在同一条边界上的任何定向曲面,旋度通量相同.

梯度没有环量,旋度没有净源,
\begin{equation}
  \nabla\times(\nabla u)=\mathbf{0},
  \qquad
  \nabla\cdot(\nabla\times\mathbf{A})=0.
\end{equation}
混合偏导可交换时,直角分量可以直接看出.完整的指标证明在后面.

\begin{example}
$\mathbf{A}=-y\,\mathbf{e}_x+x\,\mathbf{e}_y-z\,\mathbf{e}_z$,$L$ 是 $xy$ 平面上半径为 $R$ 的圆.
\end{example}
\begin{solution}
$\nabla\times\mathbf{A}=2\mathbf{e}_z$.圆上 $x=R\cos\phi$,$y=R\sin\phi$,
\begin{equation}
  \mathbf{A}\cdot\mathrm{d}\boldsymbol{\ell}=R^2\,\mathrm{d}\phi,
  \qquad
  \oint_L\mathbf{A}\cdot\mathrm{d}\boldsymbol{\ell}=2\pi R^2.
\end{equation}
圆盘上 $\int_S 2\,\mathrm{d}S=2\cdot\pi R^2$,与线积分一致.
\end{solution}

\begin{example}
无限长直导线的磁场 $\mathbf{B}=(\mu_0 I/(2\pi\rho))\,\mathbf{e}_{\phi}$ 在 $\rho>0$ 处无旋.
\end{example}
\begin{solution}
只有 $\phi$ 分量,$h_{\phi} B_{\phi}=\mu_0 I/(2\pi)$,是常数.由\eqref{eq:curl-orthogonal},旋度的每个分量都是这个常数对某个坐标求导,结果是零.环路如果把导线围在里面,环量是 $\mu_0 I$,旋度的奇性集中在 $\rho=0$.
\end{solution}

\subsection{拉普拉斯算子}
\idx{拉普拉斯算子}{lapulasisuanzi}

$\nabla^2 u=\nabla\cdot(\nabla u)$.把\eqref{eq:grad-orthogonal}的分量代进\eqref{eq:div-orthogonal},
\begin{equation}
  \nabla^2 u
  =\frac{1}{h_1 h_2 h_3}
    \sum_{i=1}^{3}
    \frac{\partial}{\partial q_i}
    \Biggl(\frac{h_1 h_2 h_3}{h_i^2}\frac{\partial u}{\partial q_i}\Biggr).
  \label{eq:laplace-orthogonal}
\end{equation}
不能把直角坐标的二阶偏导照抄过来,再把 $\mathrm{d}x$ 换成 $h_i\,\mathrm{d}q_i$.那样会漏掉尺度因子的导数.直角坐标里三个尺度因子都是 $1$,上式就是 $\partial^2 u/\partial x^2+\partial^2 u/\partial y^2+\partial^2 u/\partial z^2$.

柱坐标代入 $h_{\rho}=1$, $h_{\phi}=\rho$, $h_z=1$,
\begin{equation}
  \nabla^2 u
  =\frac{1}{\rho}\frac{\partial}{\partial\rho}
      \Bigl(\rho\frac{\partial u}{\partial\rho}\Bigr)
   +\frac{1}{\rho^2}\frac{\partial^2 u}{\partial\phi^2}
   +\frac{\partial^2 u}{\partial z^2}.
  \label{eq:laplace-cyl}
\end{equation}
第一项展开是 $\partial^2 u/\partial\rho^2+(1/\rho)\partial u/\partial\rho$.多出来的 $(1/\rho)\partial u/\partial\rho$ 就是把二阶偏导直接换元时所丢掉的项.

球坐标代入 $h_r=1$, $h_{\theta}=r$, $h_{\phi}=r\sin\theta$,
\begin{equation}
  \nabla^2 u
  =\frac{1}{r^2}\frac{\partial}{\partial r}
      \Bigl(r^2\frac{\partial u}{\partial r}\Bigr)
   +\frac{1}{r^2\sin\theta}\frac{\partial}{\partial\theta}
      \Bigl(\sin\theta\frac{\partial u}{\partial\theta}\Bigr)
   +\frac{1}{r^2\sin^2\theta}\frac{\partial^2 u}{\partial\phi^2}.
  \label{eq:laplace-sph}
\end{equation}

\begin{example}
检验 $\ln\rho$ 和 $1/r$.
\end{example}
\begin{solution}
$u=\ln\rho$ 只依赖 $\rho$.由\eqref{eq:laplace-cyl},
\begin{equation}
  \nabla^2\ln\rho
  =\frac{1}{\rho}\frac{\partial}{\partial\rho}
    \Bigl(\rho\cdot\frac{1}{\rho}\Bigr)
  =0,
  \qquad\rho>0.
\end{equation}
若误用 $\partial^2(\ln\rho)/\partial\rho^2=-1/\rho^2$,结果不是零.二维里这个函数的奇性在原点,$\nabla^2\ln\rho=2\pi\delta^2(\boldsymbol{\rho})$.

$u=1/r$ 只依赖 $r$.$r^2\partial u/\partial r=-1$,再对 $r$ 求导是零,所以 $\nabla^2(1/r)=0$（$r>0$）.它是球外拉普拉斯方程的基本解.
\end{solution}

\section{指标记号}
\idx{指标记号}{zhibiaojihao}
\idx{克罗内克符号}{keluoneike}
\idx{列维-奇维塔符号}{lieweiqiweita}

\subsection{求和约定与两个符号}
直角坐标里指标 $i,j,k$ 各取 $1,2,3$.同一个单项里重复出现的指标要求和,求和号不再写出.自由指标在等式两边必须相同,并且可以改名;哑指标只出现在求和里,也可以改名.

克罗内克符号 $\delta_{ij}$ 在 $i=j$ 时是 $1$,否则是 $0$.因此 $\delta_{ii}=3$,$A_i\delta_{ij}=A_j$.

列维--奇维塔符号 $\varepsilon_{ijk}$ 在 $(i,j,k)$ 是 $(1,2,3)$ 的偶置换时为 $+1$,奇置换时为 $-1$,有重复指标时为 $0$.它完全反对称.叉积和旋度就是
\begin{equation}
  (\mathbf{A}\times\mathbf{B})_i=\varepsilon_{ijk}A_j B_k,
  \qquad
  (\nabla\times\mathbf{A})_i=\varepsilon_{ijk}\partial_j A_k.
\end{equation}

最常用的缩并是
\begin{equation}
  \varepsilon_{ijk}\varepsilon_{ilm}
  =\delta_{jl}\delta_{km}-\delta_{jm}\delta_{kl}.
  \label{eq:eps-delta}
\end{equation}
两边都对 $j,k$ 和 $l,m$ 反对称,而且在 $(j,k)=(l,m)=(1,2)$ 时两边都是 $1$,所以作为张量相等.再缩并一次,
\begin{equation}
  \varepsilon_{ijk}\varepsilon_{ijm}=2\delta_{km},
  \qquad
  \varepsilon_{ijk}\varepsilon_{ijk}=6.
\end{equation}

\begin{example}
用\eqref{eq:eps-delta} 证明矢量三重积.
\end{example}
\begin{solution}
\begin{align}
  \bigl(\mathbf{A}\times(\mathbf{B}\times\mathbf{C})\bigr)_i
  &=\varepsilon_{ijk}A_j(\mathbf{B}\times\mathbf{C})_k
   =\varepsilon_{ijk}\varepsilon_{klm}A_j B_l C_m \\
  &=(\delta_{il}\delta_{jm}-\delta_{im}\delta_{jl})A_j B_l C_m \\
  &=A_j C_j B_i-A_j B_j C_i.
\end{align}
这就是 $(\mathbf{A}\cdot\mathbf{C})\mathbf{B}-(\mathbf{A}\cdot\mathbf{B})\mathbf{C}$.拉格朗日恒等式\eqref{eq:lagrange-identity} 是同一缩并再点上另一个叉积.
\end{solution}

\subsection{导数恒等式}
偏导与 $\varepsilon,\delta$ 可交换.因此
\begin{equation}
  \bigl(\nabla\times(\nabla f)\bigr)_i
  =\varepsilon_{ijk}\partial_j\partial_k f=0,
\end{equation}
因为 $\partial_j\partial_k f$ 对 $j,k$ 对称,而 $\varepsilon_{ijk}$ 反对称.同样
\begin{equation}
  \nabla\cdot(\nabla\times\mathbf{A})
  =\varepsilon_{ijk}\partial_i\partial_j A_k=0.
\end{equation}
散度与旋度的复合给出
\begin{equation}
  \nabla\times(\nabla\times\mathbf{A})
  =\nabla(\nabla\cdot\mathbf{A})-\nabla^2\mathbf{A}.
  \label{eq:curl-curl}
\end{equation}
左边的第 $i$ 分量是 $\varepsilon_{ijk}\partial_j(\varepsilon_{klm}\partial_l A_m)$,用\eqref{eq:eps-delta} 拆开就是 $\partial_i(\partial_m A_m)-\partial_m\partial_m A_i$.

乘积法则按分量求导即可:
\begin{align}
  \nabla\cdot(f\mathbf{A})
  &=f\,\nabla\cdot\mathbf{A}+\mathbf{A}\cdot\nabla f, \\
  \nabla\times(f\mathbf{A})
  &=f\,\nabla\times\mathbf{A}+\nabla f\times\mathbf{A}, \\
  \nabla\cdot(\mathbf{A}\times\mathbf{B})
  &=\mathbf{B}\cdot(\nabla\times\mathbf{A})
    -\mathbf{A}\cdot(\nabla\times\mathbf{B}).
\end{align}
同一缩并还给出
\begin{equation}
\begin{aligned}
  \mathbf{u}\times(\nabla\times\mathbf{u})
  &=\nabla\Bigl(\frac12|\mathbf{u}|^2\Bigr)
   -(\mathbf{u}\cdot\nabla)\mathbf{u}, \\
  \nabla\times(\mathbf{u}\times\mathbf{v})
  &=(\mathbf{v}\cdot\nabla)\mathbf{u}
   -(\mathbf{u}\cdot\nabla)\mathbf{v}
   +\mathbf{u}\,(\nabla\cdot\mathbf{v})
   -\mathbf{v}\,(\nabla\cdot\mathbf{u}).
\end{aligned}
\end{equation}
反对称与对称相乘是零,所以 $\varepsilon_{ijk}\delta_{jk}=0$.迹可以循环,$\operatorname{tr}(BC)=B_{ij}C_{ji}=\operatorname{tr}(CB)$,这里只用到分量乘法可交换.

不可压缩流动的纳维--斯托克斯方程是
\begin{equation}
  \rho\bigl(\partial_t u_i+u_j\partial_j u_i\bigr)
  =-\partial_i p+\mu\,\partial_j\partial_j u_i,
  \qquad
  \partial_t\rho+\partial_i(\rho u_i)=0.
\end{equation}
涡量 $\boldsymbol{\omega}=\nabla\times\mathbf{u}$ 的分量是 $\omega_i=\varepsilon_{ijk}\partial_j u_k$.

\begin{example}
真空中的麦克斯韦方程写成散度和旋度.
\end{example}
\begin{solution}
\begin{equation}
\begin{aligned}
  \nabla\cdot\mathbf{E}&=\frac{\rho}{\varepsilon_0}, &
  \nabla\cdot\mathbf{B}&=0, \\
  \nabla\times\mathbf{E}&=-\frac{\partial\mathbf{B}}{\partial t}, &
  \nabla\times\mathbf{B}
  &=\mu_0\mathbf{J}+\mu_0\varepsilon_0\frac{\partial\mathbf{E}}{\partial t}.
\end{aligned}
\end{equation}
对法拉第定律取散度,左边是零,所以 $\partial_t(\nabla\cdot\mathbf{B})=0$.对安培--麦克斯韦定律取散度,并用 $\nabla\cdot\mathbf{E}=\rho/\varepsilon_0$,得到连续性方程 $\partial_t\rho+\nabla\cdot\mathbf{J}=0$.对法拉第定律再取旋度,用\eqref{eq:curl-curl} 和 $\nabla\cdot\mathbf{B}=0$,在 $\mathbf{J}=\mathbf{0}$、$\rho=0$ 时得到波动方程
\begin{equation}
  \nabla^2\mathbf{B}-\mu_0\varepsilon_0\frac{\partial^2\mathbf{B}}{\partial t^2}=\mathbf{0}.
\end{equation}
\end{solution}

\subsection{行列式与正交矩阵}
二维行列式可以写成 $\det M=\varepsilon_{ij}M_{1i}M_{2j}$.三维是 $\det M=\varepsilon_{ijk}M_{1i}M_{2j}M_{3k}$.把 $MN$ 的行代入,括号中的量对三个新指标完全反对称,而且在 $(1,2,3)$ 上等于 $\det N$,所以
\begin{equation}
  \varepsilon_{ijk}N_{ai}N_{bj}N_{ck}=(\det N)\,\varepsilon_{abc},
\end{equation}
从而 $\det(MN)=\det M\det N$.

定义 $A_{jk}=\varepsilon_{jkl}a_l$.则 $(A\mathbf{b})_j=\varepsilon_{jkl}a_l b_k=(\mathbf{a}\times\mathbf{b})_j$.叉积就是这个反对称矩阵作用在 $\mathbf{b}$ 上.

\begin{example}
泡利矩阵
\begin{equation}
  \sigma^x=\begin{pmatrix}0&1\\1&0\end{pmatrix},\quad
  \sigma^y=\begin{pmatrix}0&-\mathrm{i}\\\mathrm{i}&0\end{pmatrix},\quad
  \sigma^z=\begin{pmatrix}1&0\\0&-1\end{pmatrix}.
\end{equation}
\end{example}
\begin{solution}
$(\sigma^x\sigma^y)_{ij}=\sigma^x_{ik}\sigma^y_{kj}$.对角元
\begin{equation}
  (\sigma^x\sigma^y)_{11}=\sigma^x_{12}\sigma^y_{21}=\mathrm{i},
  \qquad
  (\sigma^x\sigma^y)_{22}=\sigma^x_{21}\sigma^y_{12}=-\mathrm{i}.
\end{equation}
非对角元用到 $\sigma^y$ 的对角元,它们是零.所以 $\sigma^x\sigma^y=\mathrm{i}\,\sigma^z$.同样 $\sigma^y\sigma^x=-\mathrm{i}\,\sigma^z$.

迹 $\operatorname{tr}\sigma^x=\sigma^x_{ii}=0$,三条泡利矩阵的迹都是零,$\operatorname{tr}(\sigma^x\sigma^y)=\mathrm{i}-\mathrm{i}=0$.二维行列式
\begin{equation}
  \det\sigma^x=\varepsilon_{21}\sigma^x_{12}\sigma^x_{21}=-1,
  \qquad
  \det\sigma^y=\varepsilon_{21}\sigma^y_{12}\sigma^y_{21}=-1.
\end{equation}
因而 $\det(\sigma^x\sigma^y)=\det\sigma^x\det\sigma^y=1$,与 $\mathrm{i}\,\sigma^z$ 的行列式一致.
\end{solution}

旋转矩阵 $R^{\mathsf T}R=I$ 写成 $R_{ki}R_{kj}=\delta_{ij}$.两个矢量的点积
\begin{equation}
  (R\mathbf{a})\cdot(R\mathbf{b})=R_{ki}a_i\,R_{kj}b_j=\delta_{ij}a_i b_j=\mathbf{a}\cdot\mathbf{b}.
\end{equation}
模和夹角都被保持.

\section{张量}
\idx{张量}{zhangliang}

\subsection{阶与坐标变换}
指标的个数叫做阶,也叫秩.标量是 $0$ 阶,矢量是 $1$ 阶,$n$ 维空间里 $r$ 阶张量在计入对称性之前有 $n^r$ 个分量.对称二阶张量 $T_{ij}=T_{ji}$ 的独立分量是 $n(n+1)/2$,反对称张量 $T_{ij}=-T_{ji}$ 的对角元必须为零,独立分量是 $n(n-1)/2$.二者相加正好是 $n^2$.任意二阶张量拆成
\begin{equation}
  T_{ij}
  =\frac12(T_{ij}+T_{ji})
  +\frac12(T_{ij}-T_{ji}),
\end{equation}
第一项对称,第二项反对称.位移梯度 $\partial_i u_j$ 的对称部分是应变,反对称部分是无穷小转动.

二阶张量是把一个矢量送到另一个矢量的线性映射,$u_i=T_{ij}v_j$.三阶张量可以吃进两个矢量,$w_i=C_{ijk}u_j v_k$,对每个矢量分别线性.

直角坐标的旋转 $x'_i=R_{ij}x_j$ 下,矢量分量按
\begin{equation}
  v'_i=R_{ij}v_j
\end{equation}
变换.$2$ 阶张量的每个指标各乘一个 $R$,
\begin{equation}
  T'_{ij}=R_{ik}R_{jl}T_{kl}.
\end{equation}
若 $u_i=T_{ij}v_j$,则
\begin{equation}
  T'_{ij}v'_j
  =R_{ik}R_{jl}T_{kl}\,R_{jm}v_m
  =R_{ik}u_k,
\end{equation}
输出仍按矢量变换.$\delta_{ij}$ 在旋转下不变,因为 $R_{ik}R_{jk}=\delta_{ij}$.完全反对称的 $\varepsilon_{ijk}$ 在正常旋转下也不变,反射会多一个 $\det R=-1$.旋转矩阵本身描写的是坐标变换,不按这个规则变成张量.

应力 $\sigma_{ij}$ 把面元法向 $\mathrm{d}S_j$ 变成力 $\mathrm{d}F_i=\sigma_{ij}\,\mathrm{d}S_j$.应变描述位移梯度的对称部分.在线性弹性里,二者由弹性张量联系,$\sigma_{ij}=C_{ijkl}\varepsilon_{kl}$.各向同性时 $C$ 只用两个独立常数,拉梅常数 $\lambda,\mu$,
\begin{equation}
  \sigma_{ij}=\lambda\,\delta_{ij}\,\varepsilon_{kk}+2\mu\,\varepsilon_{ij}.
\end{equation}
电导张量把电场变成电流,$J_i=\sigma_{ij}E_j$.晶体里 $\sigma_{ij}$ 不必是 $\delta_{ij}$ 的倍数,电流可以不平行于电场.压电张量 $d_{ijk}$ 把应力变成极化,$P_i=d_{ijk}\sigma_{jk}$.

\subsection{各向同性张量}
\idx{各向同性张量}{gexiangtongxing}
旋转下形式不变的张量由 $\delta$ 和 $\varepsilon$ 拼出来.$0$ 阶就是常数.$1$ 阶只有零矢量.$2$ 阶是 $\lambda\delta_{ij}$.$3$ 阶是 $\lambda\varepsilon_{ijk}$.$4$ 阶有三项,
\begin{equation}
  T_{ijkl}
  =\alpha\,\delta_{ij}\delta_{kl}
  +\beta\,\delta_{ik}\delta_{jl}
  +\gamma\,\delta_{il}\delta_{jk}.
\end{equation}
弹性张量还要满足指标对称,所以 $\beta=\gamma$,独立常数减成两个,与上面的 $\lambda,\mu$ 对应.

\subsection{度规、联络与曲率}
\idx{度规}{dugui}
\idx{克里斯托费尔符号}{kelisituoer}
曲线坐标里,坐标差 $\mathrm{d}x^\mu$ 不是弧长.度规把弧长写成
\begin{equation}
  \mathrm{d}s^2=g_{\mu\nu}\,\mathrm{d}x^\mu\mathrm{d}x^\nu.
\end{equation}
平面极坐标 $x^1=r$,$x^2=\vartheta$,
\begin{equation}
  \mathrm{d}s^2=\mathrm{d}r^2+r^2\mathrm{d}\vartheta^2,
  \qquad
  g_{rr}=1,\quad g_{\vartheta\vartheta}=r^2,
\end{equation}
非对角元是零.$g^{rr}=1$,$g^{\vartheta\vartheta}=1/r^2$.尺度因子就是 $h_r=\sqrt{g_{rr}}=1$,$h_\vartheta=\sqrt{g_{\vartheta\vartheta}}=r$,与第~\ref{sec:coordinates}~节一致.

克里斯托费尔符号
\begin{equation}
  \Gamma^\rho{}_{\mu\nu}
  =\frac12 g^{\rho\sigma}
    \bigl(
      \partial_\mu g_{\nu\sigma}
      +\partial_\nu g_{\mu\sigma}
      -\partial_\sigma g_{\mu\nu}
    \bigr)
  \label{eq:christoffel}
\end{equation}
描写坐标基如何逐点变化.协变导数是
\begin{equation}
  \nabla_\mu V^\nu
  =\partial_\mu V^\nu+\Gamma^\nu{}_{\mu\rho}V^\rho,
  \qquad
  \nabla_\mu V_\nu
  =\partial_\mu V_\nu-\Gamma^\rho{}_{\mu\nu}V_\rho.
\end{equation}
$\Gamma$ 不是张量:它在直线坐标里可以全部为零,换到极坐标就不是零.曲率张量
\begin{equation}
\begin{aligned}
  R^\rho{}_{\sigma\mu\nu}
  &=\partial_\mu\Gamma^\rho{}_{\nu\sigma}
   -\partial_\nu\Gamma^\rho{}_{\mu\sigma} \\
  &\quad
   +\Gamma^\rho{}_{\mu\lambda}\Gamma^\lambda{}_{\nu\sigma}
   -\Gamma^\rho{}_{\nu\lambda}\Gamma^\lambda{}_{\mu\sigma}
\end{aligned}
\end{equation}
才按张量变换.缩并 $R_{\mu\nu}=R^\rho{}_{\mu\rho\nu}$,$R=g^{\mu\nu}R_{\mu\nu}$,爱因斯坦张量 $G_{\mu\nu}=R_{\mu\nu}-\frac12 R g_{\mu\nu}$ 与物质能动张量的关系是 $G_{\mu\nu}+\Lambda g_{\mu\nu}=(8\pi G/c^4)T_{\mu\nu}$.

平面极坐标里非零的只有
\begin{equation}
  \Gamma^r{}_{\vartheta\vartheta}=-r,
  \qquad
  \Gamma^\vartheta{}_{\,r\vartheta}
  =\Gamma^\vartheta{}_{\,\vartheta r}=\frac1r.
  \label{eq:polar-gamma}
\end{equation}
例如 $\Gamma^r{}_{rr}=\frac12 g^{rr}(\partial_r g_{rr})=0$.径向测地线方程是
\begin{equation}
  \frac{\mathrm{d}^2 r}{\mathrm{d}s^2}
  +\Gamma^r{}_{\vartheta\vartheta}
  \Bigl(\frac{\mathrm{d}\vartheta}{\mathrm{d}s}\Bigr)^2=0,
\end{equation}
代入\eqref{eq:polar-gamma} 得到 $\mathrm{d}^2 r/\mathrm{d}s^2=r(\mathrm{d}\vartheta/\mathrm{d}s)^2$.这是直线在极坐标里的加速度,不是平面被弯曲.

黎曼曲率的一个分量是
\begin{equation}
\begin{aligned}
  R^r{}_{\vartheta r\vartheta}
  &=\partial_r\Gamma^r{}_{\vartheta\vartheta}
   -\partial_\vartheta\Gamma^r{}_{\,r\vartheta} \\
  &\quad
   +\Gamma^r{}_{\,r\lambda}\Gamma^\lambda{}_{\vartheta\vartheta}
   -\Gamma^r{}_{\vartheta\lambda}\Gamma^\lambda{}_{\,r\vartheta}.
\end{aligned}
\end{equation}
第一项 $\partial_r(-r)=-1$.$\Gamma^r{}_{\,r\vartheta}=0$,第二、三项是零.第四项里 $\lambda=\vartheta$ 的贡献是 $(-r)(1/r)=-1$,前面再有一个减号,贡献 $+1$.合计 $-1+1=0$.

二维的曲率张量只有一个独立分量.一般地,先在 $n$ 维里选出一对反对称指标,有 $m=n(n-1)/2$ 种选法;两对之间再对称,独立个数是 $m(m+1)/2$.二维是 $1$,三维是 $6$.上面这一个分量为零,所以平面的曲率张量是零,标量曲率也是零.极坐标里出现的 $\Gamma$ 只是坐标效应.

\subsection{惯量张量}
\idx{惯量张量}{guanliangzhangliang}
角速度与角动量由惯量张量联系,$L_i=I_{ij}\omega_j$.质点 $m_a$ 位于 $\mathbf{r}^{(a)}$ 时
\begin{equation}
  I_{ij}
  =\sum_a m_a
    \bigl(
      |\mathbf{r}^{(a)}|^2\delta_{ij}
      -r^{(a)}_i r^{(a)}_j
    \bigr)
\end{equation}
$I_{ij}=I_{ji}$.存在主轴使它对角,$L_i=I_i\omega_i$,三个主转动惯量一般不同,$\mathbf{L}$ 不必平行于 $\boldsymbol{\omega}$.

\subsection{缩并的次序}
若干张量全部缩并以后得到标量.例如 $A_{ij}B_{jk}C_{ki}$ 可以先缩 $B$ 和 $C$,再与 $A$ 缩,也可以先缩 $A$ 和 $B$.求和只是有限项相加,次序可以交换,两个算法给出同一个数.指标图把每个张量画成一个结点,每条边是一个被缩并的指标;图的连通方式和最后的标量,不依赖于先算哪一个结点.

下面的题把本章的恒等式和例题再算一遍.指标题要用求和约定,不能把分量逐个展开以后只检验一个特例.

\begin{problemset}
\renewcommand{\theenumi}{\thechapter.\arabic{enumi}}
\renewcommand{\labelenumi}{\theenumi.}
\item 证明 $\delta_{ij}a_j=a_i$, $\delta_{ij}\delta_{jk}=\delta_{ik}$, $\delta_{ii}=3$.

\item 由 $\varepsilon_{ijk}$ 的定义说明交换任意两个指标变号,并由此证明 $\mathbf{u}\times\mathbf{u}=\mathbf{0}$.

\item 从 $(\mathbf{u}\times\mathbf{v})_i=\varepsilon_{ijk}u_j v_k$ 推出标量三重积 $\mathbf{u}\cdot(\mathbf{v}\times\mathbf{w})=\varepsilon_{ijk}u_i v_j w_k$.说明循环换位不变号,交换两个矢量变号.

\item 证明 $\varepsilon_{ijk}\varepsilon_{klm}=\delta_{il}\delta_{jm}-\delta_{im}\delta_{jl}$,并用它得到
\begin{equation*}
  \mathbf{u}\times(\mathbf{v}\times\mathbf{w})
  =(\mathbf{u}\cdot\mathbf{w})\mathbf{v}-(\mathbf{u}\cdot\mathbf{v})\mathbf{w}.
\end{equation*}

\item 由 $R_{ki}R_{kj}=\delta_{ij}$ 证明正交变换保持内积,从而保持长度和夹角.

\item 用指标证明 $\nabla\times\nabla\varphi=\mathbf{0}$, $\nabla\cdot(\nabla\times\mathbf{u})=0$,以及
\begin{equation*}
  \nabla\times(\nabla\times\mathbf{E})
  =\nabla(\nabla\cdot\mathbf{E})-\nabla^2\mathbf{E}.
\end{equation*}

\item 把不可压缩流动的纳维--斯托克斯方程和连续性方程写成指标形式,并写出涡量 $\omega_i=\varepsilon_{ijk}\partial_j u_k$.

\item 用余子式求
\begin{equation*}
  A=\begin{pmatrix}2&1\\1&2\end{pmatrix}
\end{equation*}
的逆,并解 $A\mathbf{x}=(1,0)^{\mathsf T}$.

\item 对称二阶张量有 $n(n+1)/2$ 个独立分量,反对称二阶张量有 $n(n-1)/2$ 个.证明二者之和是 $n^2$,并且任意 $T_{ij}$ 可以唯一地拆成这两部分.

\item 设 $T'_{ij}=R_{ik}R_{jl}T_{kl}$, $v'_j=R_{jm}v_m$.验证 $u_i=T_{ij}v_j$ 仍按矢量变换.

\item 各向同性电导 $\sigma_{ij}=\sigma\delta_{ij}$.由 $J_i=\sigma_{ij}E_j$ 说明 $\mathbf{J}$ 平行于 $\mathbf{E}$.举一个非各向同性的 $2\times 2$ 例子,使电流不平行于电场.

\item 极坐标 $\mathrm{d}s^2=\mathrm{d}r^2+r^2\mathrm{d}\vartheta^2$.算出非零的克里斯托费尔符号,并验证 $R^r{}_{\vartheta r\vartheta}=0$.说明这个例子为什么表明 $\Gamma$ 不是张量.

\item 柱坐标里验证 $h_\rho=1$, $h_\phi=\rho$, $h_z=1$,并计算 $\nabla\cdot\mathbf{e}_\rho$ 和 $\nabla^2\ln\rho$ ($\rho>0$).指出若把拉普拉斯算子写成 $\partial^2/\partial\rho^2$,第二问会错在哪里.

\item 用泡利矩阵的指标乘法计算 $\sigma^x\sigma^y$ 和 $\sigma^y\sigma^x$ 的对角元,以及 $\det\sigma^x$, $\det\sigma^y$ 和 $\det(\sigma^x\sigma^y)$.
\end{problemset}

